\documentclass[a4paper,11pt]{article}

% ── Packages ──────────────────────────────────────────────────────────────────
\usepackage[T1]{fontenc}
\usepackage[utf8]{inputenc}
\usepackage[english]{babel}

\usepackage{geometry}
\geometry{margin=2.5cm}

\usepackage{amsmath, amssymb}
\usepackage{booktabs}
\usepackage{tabularx}
\usepackage{multirow}
\usepackage{graphicx}
\usepackage{float}
\usepackage{caption}
\usepackage{subcaption}
\usepackage{xcolor}
\usepackage{listings}
\usepackage{hyperref}
\usepackage{enumitem}
\usepackage{parskip}
\usepackage{titlesec}
\usepackage{fancyhdr}

% ── Code listing style ─────────────────────────────────────────────────────────
\definecolor{codebg}{RGB}{245,245,245}
\definecolor{codegreen}{RGB}{0,128,0}
\definecolor{codepurple}{RGB}{128,0,128}
\definecolor{codered}{RGB}{180,0,0}

\lstdefinestyle{pythonstyle}{
  backgroundcolor=\color{codebg},
  commentstyle=\color{codegreen},
  keywordstyle=\color{codepurple}\bfseries,
  stringstyle=\color{codered},
  basicstyle=\ttfamily\footnotesize,
  breaklines=true,
  captionpos=b,
  keepspaces=true,
  numbers=left,
  numberstyle=\tiny\color{gray},
  showspaces=false,
  showstringspaces=false,
  language=Python,
  frame=single,
  rulecolor=\color{gray!40},
  xleftmargin=1.5em,
  framexleftmargin=1.5em,
}
\lstset{style=pythonstyle}

% ── Header / footer ────────────────────────────────────────────────────────────
\pagestyle{fancy}
\fancyhf{}
\rhead{\small INFO-H-515 — Project 2025–2026}
\lhead{\small Distributed RAG System}
\cfoot{\thepage}
\renewcommand{\headrulewidth}{0.4pt}

% ── Section title formatting ───────────────────────────────────────────────────
\titleformat{\section}{\large\bfseries}{}{0em}{}[\titlerule]
\titleformat{\subsection}{\normalsize\bfseries}{}{0em}{}
\titleformat{\subsubsection}{\normalsize\itshape}{}{0em}{}

% ── Hyperref ───────────────────────────────────────────────────────────────────
\hypersetup{
  colorlinks=true,
  linkcolor=blue!60!black,
  citecolor=blue!60!black,
  urlcolor=blue!60!black,
}

% ══════════════════════════════════════════════════════════════════════════════
\begin{document}

% ── Title page ────────────────────────────────────────────────────────────────
\begin{titlepage}
  \centering
  \vspace*{2cm}
  {\Large \textsc{Université Libre de Bruxelles}\\[0.3em]
   \textsc{École Polytechnique de Bruxelles}\\[0.3em]
   \textsc{Faculté des Sciences}\par}
  \vspace{1.5cm}
  \rule{\linewidth}{0.5pt}\\[0.4cm]
  {\LARGE \textbf{INFO-H-515: Big Data}\\[0.3em]
   Distributed Data Management and Scalable Analytics\\[0.3em]
   \large Project Report 2025–2026\par}
  \rule{\linewidth}{0.5pt}\\[1.5cm]

  {\large \textbf{Distributed Retrieval-Augmented Generation (RAG) System\\
  for Automatic Exam Question Generation}}\\[2cm]

  \begin{tabular}{ll}
    \textbf{Authors:} & Author One \\
                      & Author Two \\
                      & Author Three \\
                      & Author Four \\[1em]
    \textbf{Academic year:} & 2025–2026 \\
    \textbf{Date:}          & May 2026 \\
  \end{tabular}

  \vfill
  {\small Supervisors: Dimitris Sacharidis, Gianluca Bontempi, Cédric Simar,
  Martin Colot, Antonios Kontaxakis}
\end{titlepage}

% ── Table of contents ─────────────────────────────────────────────────────────
\tableofcontents
\newpage

% ══════════════════════════════════════════════════════════════════════════════
\section{Introduction}

The goal of this project is to design and implement a \textbf{Distributed
Retrieval-Augmented Generation (RAG) system} capable of automatically
generating exam questions from a corpus of course-material PDFs.
The system is built on top of Apache Spark (PySpark) to ensure scalability,
and integrates classical and neural embedding techniques alongside a Large
Language Model (LLM) for question generation.

The project is divided into three tasks:
\begin{enumerate}[noitemsep]
  \item \textbf{Task 1} — Ingestion, preprocessing, chunking, embedding, and
        storage of the PDF corpus.
  \item \textbf{Task 2} — Semantic retrieval of relevant chunks and LLM-based
        exam question generation.
  \item \textbf{Task 3} — Comprehensive evaluation of retrieval quality,
        generation quality, and system performance.
\end{enumerate}

This report describes the architecture, design choices, and experimental
results for each task.

% ══════════════════════════════════════════════════════════════════════════════
\section{Overall System Architecture}

\subsection{Pipeline Overview}

The pipeline follows a classic RAG architecture adapted to a distributed
setting, as depicted below:

\begin{enumerate}[noitemsep]
  \item PDF documents are read in parallel from an input directory.
  \item Each document is converted to page-level raw text, cleaned, and
        split into overlapping word-based chunks.
  \item Each chunk is embedded using one of three strategies: TF-IDF,
        Word2Vec, or SBERT.
  \item Chunks and their embeddings are stored in Parquet format on HDFS,
        alongside rich traceability metadata.
  \item At query time, the user query is embedded with the same strategy,
        and the top-$k$ most similar chunks are retrieved via a distance
        metric.
  \item Retrieved chunks are injected into a structured prompt, which is
        sent to a local LLM to generate grounded exam questions with
        citations.
\end{enumerate}

\subsection{Distributed Processing with PySpark}

The pipeline makes extensive use of PySpark's RDD API:
\begin{itemize}[noitemsep]
  \item \texttt{sc.parallelize} distributes the list of PDF file paths across
        available executor cores.
  \item \texttt{flatMap} is used to expand each PDF into its constituent pages,
        and again to expand each cleaned page into its chunks.
  \item \texttt{map} applies the cleaning function page by page in a
        data-parallel fashion.
  \item \texttt{filter} (optionally enabled) removes empty or near-empty
        pages before chunking.
  \item \texttt{collect} gathers chunks onto the driver before fitting
        global models such as TF-IDF or SBERT, which require access to the
        full corpus.
\end{itemize}

The architecture is horizontally scalable: adding executor cores or nodes
directly reduces the ingestion and cleaning time, as each PDF page is an
independent unit of work. The main bottleneck is the embedding step, which
is CPU/GPU-bound and is discussed further in Task~3.

% ══════════════════════════════════════════════════════════════════════════════
\section{Task 1: Preprocessing, Tokenization, and Embeddings}

\subsection{Configuration}

All pipeline parameters are centralised in a single \texttt{RAGConfig}
dataclass, making the notebook reproducible and easy to reconfigure:

\begin{lstlisting}
@dataclass
class RAGConfig:
    pdf_path: str = "course_material"
    chunk_size_words: int  = 200
    chunk_overlap_words: int = 50
    top_k: int = 4
    similarity_metric: str = "cosine"   # {"cosine", "euclidean"}
    embedding_strategy: str = "sbert"   # {"tfidf", "sbert", "word2vec"}
    sbert_model_name: str = "sentence-transformers/all-MiniLM-L6-v2"
    gen_model_name: str = "google/flan-t5-small"
    max_new_tokens: int = 256
\end{lstlisting}

\subsection{PDF Ingestion}

PDF files are discovered with \texttt{Path.glob("*.pdf")} and the resulting
list of paths is distributed via \texttt{sc.parallelize}. Text extraction
relies on the \texttt{pypdf} library (\texttt{PdfReader}), which provides
page-level access without requiring a native PDF renderer. For each page,
the following metadata fields are recorded immediately at extraction time:

\begin{table}[H]
  \centering
  \caption{Metadata fields captured at ingestion}
  \begin{tabularx}{\linewidth}{lX}
    \toprule
    \textbf{Field} & \textbf{Description} \\
    \midrule
    \texttt{pdf\_name}  & Filename of the source PDF \\
    \texttt{page\_num}  & 1-based page index within the document \\
    \texttt{text}       & Raw extracted text for that page \\
    \bottomrule
  \end{tabularx}
\end{table}

\subsection{Text Cleaning}

Raw text extracted from PDFs frequently contains artefacts that degrade
retrieval quality. The \texttt{clean\_text} function applies the following
transformations in order:

\begin{enumerate}[noitemsep]
  \item \textbf{Hyphenation removal}: end-of-line hyphens (\texttt{-\textbackslash n})
        are deleted, rejoining words split across lines.
  \item \textbf{Line-break normalisation}: remaining newlines are replaced by
        spaces.
  \item \textbf{Whitespace collapsing}: runs of whitespace are reduced to a
        single space via the regex \texttt{\textbackslash s+}.
  \item \textbf{Strip}: leading and trailing whitespace is removed.
\end{enumerate}

A filter step (currently commented out but configurable) discards pages whose
cleaned text is shorter than 50 characters, eliminating blank or
nearly-blank pages that would otherwise introduce noise into the index.

All cleaning operations are applied with \texttt{pages\_rdd.map(clean\_text)},
keeping the transformation fully distributed.

\subsection{Chunking}

\subsubsection{Rationale for chunk size and overlap}

The chunk size is set to \textbf{200 words} with an overlap of \textbf{50 words}
(25\%), giving a step size of 150 words. This configuration was chosen for
the following reasons:

\begin{itemize}[noitemsep]
  \item 200 words corresponds roughly to one to two paragraphs of academic
        text, a natural semantic unit that tends to cover a single concept or
        definition without excessive dilution.
  \item An overlap of 25\% ensures that definitions or formulas that straddle
        a boundary are captured in at least one chunk, reducing the risk of
        retrieving a chunk that contains only half of a key idea.
  \item The chosen size fits comfortably within the context window of small
        LLMs (e.g.\ \texttt{flan-t5-small}), and several chunks can be
        concatenated for larger models without hitting the token limit.
\end{itemize}

Both \texttt{chunk\_size\_words} and \texttt{chunk\_overlap\_words} are
exposed as configuration parameters and can be tuned without modifying the
pipeline code.

\subsubsection{Implementation}

The \texttt{word\_based\_chunk\_page} function splits the cleaned text of a
single page on whitespace and slides a window of size
\texttt{chunk\_size\_words} with step \texttt{chunk\_size\_words - chunk\_overlap\_words}.
Each chunk is assigned the following metadata:

\begin{table}[H]
  \centering
  \caption{Metadata fields added at chunking}
  \begin{tabularx}{\linewidth}{lX}
    \toprule
    \textbf{Field} & \textbf{Description} \\
    \midrule
    \texttt{chunk\_id}  & Unique identifier of the form \texttt{p\{page\}\_c\{index\}} \\
    \texttt{pdf\_name}  & Propagated from the page record \\
    \texttt{page\_num}  & Propagated from the page record \\
    \texttt{text}       & Space-joined words of the chunk \\
    \bottomrule
  \end{tabularx}
\end{table}

A guard raises \texttt{ValueError} if \texttt{overlap\_words $\geq$ chunk\_size\_words},
preventing an infinite loop. The transformation is applied with
\texttt{clean\_rdd.flatMap(word\_based\_chunk\_page)}, producing one RDD
element per chunk across all pages and all documents.

\subsection{Embeddings}

After chunking, \texttt{clean\_chunk\_rdd.collect()} brings all chunks to the
Spark driver. This is necessary because TF-IDF requires a global vocabulary
fit, Word2Vec requires a global training pass, and SBERT's batch encoding is
most efficient when run on the full corpus at once. Three embedding strategies
are supported.

\subsubsection{TF-IDF (Sparse Lexical Baseline)}

\texttt{TfidfVectorizer} from scikit-learn is fitted on all chunk texts and
produces a sparse matrix $X \in \mathbb{R}^{N \times |V|}$, where $N$ is the
number of chunks and $|V|$ the vocabulary size. TF-IDF is parameter-free in
our configuration and serves as the baseline for exact term matching.

\begin{lstlisting}
vectorizer = TfidfVectorizer()
X = vectorizer.fit_transform([c["text"] for c in chunks])
\end{lstlisting}

\subsubsection{Word2Vec (Distributional Semantics)}

A Word2Vec model is trained from scratch on the tokenised chunks using the
gensim library. Each chunk is embedded by averaging the vectors of its
known words:

\[
  \mathbf{e}_{\text{chunk}} =
  \frac{1}{|\mathcal{W}|} \sum_{w \in \mathcal{W}} \mathbf{v}_w,
  \quad \mathcal{W} = \{w \in \text{chunk} \mid w \in \text{vocab}\}
\]

If no word in the chunk is in the vocabulary (e.g.\ a chunk consisting
entirely of rare tokens), a zero vector is returned. The model uses
\texttt{vector\_size=100}, \texttt{window=5}, \texttt{min\_count=2}, and
\texttt{epochs=10}.

\subsubsection{SBERT (Contextual Dense Embeddings)}

Sentence-BERT produces contextualised embeddings via a transformer encoder
with mean pooling. The model \texttt{all-MiniLM-L6-v2} (384 dimensions) is
used for its strong quality-to-speed trade-off. Embeddings are
$\ell_2$-normalised at encoding time (\texttt{normalize\_embeddings=True}),
so that cosine similarity reduces to a dot product at retrieval time:

\begin{lstlisting}
model = SentenceTransformer(model_name)
emb = model.encode(texts, show_progress_bar=True,
                   normalize_embeddings=True)
emb = np.asarray(emb, dtype=np.float32)
\end{lstlisting}

\subsection{Storage (HDFS Parquet)}

A SparkSession is created from the existing SparkContext to gain access to the
DataFrame API required for Parquet serialisation. The
\texttt{save\_to\_hdfs\_parquet} function assembles a list of tuples
matching the following schema and writes it to HDFS:

\begin{table}[H]
  \centering
  \caption{Parquet schema for the embedding index}
  \begin{tabularx}{\linewidth}{llX}
    \toprule
    \textbf{Field} & \textbf{Spark type} & \textbf{Description} \\
    \midrule
    \texttt{id}                 & \texttt{LongType}              & Sequential integer identifier \\
    \texttt{source\_pdf}        & \texttt{StringType}            & Source PDF filename \\
    \texttt{page\_num}          & \texttt{IntegerType}           & Page number (1-based) \\
    \texttt{chunk\_id}          & \texttt{StringType}            & Chunk identifier \\
    \texttt{start\_word}        & \texttt{StringType}            & First word of the chunk \\
    \texttt{end\_word}          & \texttt{StringType}            & Last word of the chunk \\
    \texttt{chunk\_text}        & \texttt{StringType}            & Full chunk text \\
    \texttt{embedding}          & \texttt{ArrayType(FloatType)}  & Dense or sparse embedding vector \\
    \texttt{embedding\_strategy}& \texttt{StringType}            & Strategy used (\texttt{tfidf}, \texttt{word2vec}, \texttt{sbert}) \\
    \bottomrule
  \end{tabularx}
\end{table}

The DataFrame is written in \texttt{overwrite} mode to
\texttt{hdfs://namenode:8020/<output\_dir>}, which allows idempotent
re-runs. Separate output directories are used per embedding strategy.

% ══════════════════════════════════════════════════════════════════════════════
\section{Task 2: Retrieval and Generation}
\label{sec:task2}

\textit{This section will be completed once Task~2 is implemented.
The following subsections outline the planned architecture.}

\subsection{Retrieval}

Given a user query, the system will:
\begin{enumerate}[noitemsep]
  \item Embed the query using the \textbf{same} strategy and model that was
        used to build the index (ensuring comparable vector spaces).
  \item Load the Parquet index from HDFS.
  \item Compute pairwise similarity or distance between the query vector and
        all stored embeddings.
  \item Return the top-$k$ chunks ranked by score.
\end{enumerate}

The system supports two distance metrics (implemented from scratch without
external libraries, as required):

\paragraph{Cosine similarity.}
\[
  \text{sim}(\mathbf{a}, \mathbf{b})
  = \frac{\mathbf{a} \cdot \mathbf{b}}{\|\mathbf{a}\| \, \|\mathbf{b}\|}
\]
Higher values indicate greater similarity; chunks are sorted in descending
order.

\paragraph{Euclidean distance.}
\[
  d(\mathbf{a}, \mathbf{b})
  = \sqrt{\sum_{i=1}^{n} (a_i - b_i)^2}
\]
Lower values indicate greater proximity; chunks are sorted in ascending order.

\subsection{LLM Integration and Prompt Construction}

Retrieved chunks will be injected into a structured prompt that instructs the
model to generate grounded exam questions. The prompt enforces the following
constraints:
\begin{itemize}[noitemsep]
  \item Use \textbf{only} the provided sources.
  \item Cite every question and answer key with \texttt{[n]} references.
  \item Include both conceptual questions and, where relevant, formula-based
        derivations.
\end{itemize}

The generator backend is Ollama with \texttt{llama3.2:3b}, called from the
Spark driver after \texttt{.collect()} to avoid executor-side HTTP calls.

% ══════════════════════════════════════════════════════════════════════════════
\section{Task 3: Evaluation}
\label{sec:task3}

\textit{This section will be completed once Task~3 is implemented.
The subsections below describe the planned evaluation methodology.}

\subsection{Retrieval Metrics}

Precision@$k$ and Recall@$k$ will be reported for each embedding strategy
and distance metric against a manually annotated gold set of query–page
pairs.

\[
  \text{Precision@}k = \frac{|\text{retrieved}_{k} \cap \text{relevant}|}{k},
  \qquad
  \text{Recall@}k = \frac{|\text{retrieved}_{k} \cap \text{relevant}|}{|\text{relevant}|}
\]

\subsection{Generation and Grounding Metrics}

Generated questions will be evaluated on relevance, correctness,
faithfulness, citation support, clarity, and format adherence, using a
5-point rubric applied by manual inspection.

\subsection{System Performance Metrics}

Wall-clock latency, throughput (chunks/second), token counts (input and
output), and memory footprint will be reported for varying values of
\texttt{parallelism}, \texttt{chunk\_size}, and \texttt{chunk\_overlap}.
Spark UI screenshots will be included.

% ══════════════════════════════════════════════════════════════════════════════
\section{Conclusion}

Task~1 of the distributed RAG pipeline has been implemented and validated.
The PySpark notebook reads a directory of PDF files, cleans and chunks the
text in a fully distributed fashion, produces three flavours of embeddings
(TF-IDF, Word2Vec, SBERT), and stores the resulting index in HDFS Parquet
with rich traceability metadata. The configuration is centralised in
\texttt{RAGConfig}, making the system reproducible and easy to extend for
Tasks~2 and~3.

% ══════════════════════════════════════════════════════════════════════════════
\appendix
\section{Software Dependencies}

\begin{table}[H]
  \centering
  \caption{Main libraries used in the project}
  \begin{tabularx}{\linewidth}{llX}
    \toprule
    \textbf{Library} & \textbf{Version} & \textbf{Role} \\
    \midrule
    PySpark              & 3.x    & Distributed RDD / DataFrame processing \\
    pypdf                & $\geq$3 & PDF text extraction \\
    scikit-learn         & $\geq$1 & TF-IDF vectorisation \\
    gensim               & $\geq$4 & Word2Vec training \\
    sentence-transformers& $\geq$2 & SBERT encoding \\
    numpy                & $\geq$1.24 & Dense matrix operations \\
    pandas               & $\geq$2 & Evaluation DataFrames \\
    requests             & —      & Ollama HTTP client \\
    \bottomrule
  \end{tabularx}
\end{table}

\end{document}
